\begin{center}
\begin{tikzpicture}[
    thick,
    % Stile per le linee dei fermioni (linea con freccia a metà)
    fermion/.style={draw=black, postaction={decorate},
        decoration={markings, mark=at position 0.55 with {\arrow{latex}}}},
    % Stile per i bosoni W (linea ondulata)
    boson/.style={draw=black, decorate,
        decoration={snake, amplitude=1mm, segment length=3.5mm}},
    % Stile per la massa di Majorana (la X al centro)
    cross/.style={cross out, draw=black, thick, minimum size=6pt, inner sep=0pt},
    % Stile per i vertici dei nucleoni (i cerchi grigi)
    blob/.style={circle, draw=black, fill=gray!30, minimum size=6mm, inner sep=0pt}
]

% ---------------- NUCLEONE SUPERIORE ----------------
\draw[fermion] (-4.5, 2) node[left, font=\Large] {n} -- (-3, 2);
\draw[fermion] (3, 2) -- (4.5, 2) node[right, font=\Large] {p};

\draw[fermion] (-3, 2.15) -- (3, 2.15);
\draw[fermion] (-3, 2.0) -- (3, 2.0);
\draw[fermion] (-3, 1.85) -- (3, 1.85);

% ---------------- NUCLEONE INFERIORE ----------------
\draw[fermion] (-4.5, -2) node[left, font=\Large] {n} -- (-3, -2);
\draw[fermion] (3, -2) -- (4.5, -2) node[right, font=\Large] {p};

\draw[fermion] (-3, -1.85) -- (3, -1.85);
\draw[fermion] (-3, -2.0) -- (3, -2.0);
\draw[fermion] (-3, -2.15) -- (3, -2.15);

% --- DISEGNO I BLOB ---
\node[blob] at (-3, 2) {};
\node[blob] at (3, 2) {};
\node[blob] at (-3, -2) {};
\node[blob] at (3, -2) {};

% ---------------- SCAMBIO INTERNO ----------------
% Bosoni W
\draw[boson] (-0.5, 1.85) -- (0.5, 0.7);
\draw[boson] (-0.5, -1.85) -- (0.5, -0.7);

% Neutrino di Majorana (dritto, senza freccia, con la X)
\draw (0.5, 0.7) -- (0.5, -0.7);
\node[cross, label={[font=\Large]right:{$\nu_M$}}] at (0.5, 0) {};

% Elettroni uscenti
\draw[fermion] (0.5, 0.7) -- (3.5, 1.3) node[right, font=\Large] {$e^-$};
\draw[fermion] (0.5, -0.7) -- (3.5, -1.3) node[right, font=\Large] {$e^-$};

\end{tikzpicture}
\end{center}